% ===========================================================
% Глава 7: Аппроксимация функций ценности и глубокое обучение с подкреплением
% ===========================================================
\chapter[Аппроксимация функций ценности]{Аппроксимация функций ценности и глубокое обучение с подкреплением}
\label{ch:value-approx}

\begin{quote}
\itshape
Сила нейронных сетей не в том, что они запоминают; она в том, что они обобщают.
В обучении с подкреплением это различие разделяет методы, которые масштабируются, от тех, которые не масштабируются.
\end{quote}

\bigskip

Предыдущие шесть глав создали строгую основу для обучения с подкреплением, используя табличные представления.  В главе~\ref{ch:mdp} мы формализовали фреймворк марковских процессов принятия решений и вывели уравнения Беллмана; в главе~\ref{ch:dp} показали, как динамическое программирование решает их точно при наличии модели; в главах~\ref{ch:mc} и \ref{ch:td} разработали методы Монте-Карло и временной разности, которые обучаются из опыта без модели.  Эти алгоритмы обладают хорошо понятными гарантиями сходимости и чёткой математической структурой.  Однако они все разделяют одно критическое допущение: что мы можем хранить и обновлять отдельное численное значение для каждого состояния\,---\,или каждой пары состояние--действие\,---\,в задаче.

Это допущение нарушается почти мгновенно в любом реалистичном приложении.  Игра Atari представляет примерно $10^{120\,000}$ различных конфигураций экрана; роботическая рука, двигающаяся в непрерывном пространстве суставов, имеет несчётно бесконечное пространство состояний; языковая модель должна присваивать ценности каждой возможной последовательности токенов.  Ни один компьютер, ни существующий, ни будущий, не может хранить таблицу такого размера.  Даже если бы хранилище было неограниченным, обобщающие свойства табличных методов были бы плохими: обучение тому, что одно игровое состояние опасно, не даёт никакой информации о визуально схожем состоянии.

Решение --- \emph{аппроксимация функций}\index{function approximation}: вместо хранения одного числа на состояние мы обучаем параметризованную функцию $V_\mathbf{w}(s) \approx V^\pi(s)$, чьё число параметров $d$ значительно меньше $|\cS|$.  Один набор параметров одновременно кодирует информацию обо всех состояниях, и обновление параметров после посещения одного состояния автоматически изменяет предсказанные ценности близких состояний.  Это механизм обобщения.

Эта глава систематически развивает аппроксимацию функций --- от первых принципов до архитектур глубоких нейронных сетей, лежащих в основе современного обучения с подкреплением.  Раздел~\ref{sec:tabular-fail} количественно описывает проблему масштабируемости.  Раздел~\ref{sec:linear-vfa} вводит линейную аппроксимацию функций --- простейшую постановку, в которой сходимость можно доказать строго.  Раздел~\ref{sec:projected-bellman} анализирует проецированное уравнение Беллмана, характеризующее неподвижную точку TD-методов в пространстве аппроксимации функций.  Раздел~\ref{sec:semi-gradient} исследует методы полуградиента и объясняет, почему они не являются истинными методами градиентного спуска.  Раздел~\ref{sec:deadly-triad} рассматривает \emph{смертельную триаду}: сочетание аппроксимации функций, бутстрэппинга и данных с другой стратегии, способное вызвать расходимость даже при линейных аппроксиматорах.  Раздел~\ref{sec:dqn} представляет Deep Q-Networks (DQN) --- первый алгоритм, продемонстрировавший производительность уровня человека во многих разнообразных задачах, стабилизировав глубокое ОП двумя ключевыми инженерными новшествами.  Раздел~\ref{sec:dqn-extensions} рассматривает важные алгоритмические улучшения, построенные на фундаменте DQN.  Раздел~\ref{sec:convergence-vfa} обозревает, что теория сходимости может и чего не может сказать о нелинейной аппроксимации функций.  Наконец, раздел~\ref{sec:applications-vfa} обсуждает практические соображения, соединяющие теорию и инженерию.

% ===========================================================
\section{Почему табличные методы терпят неудачу}
\label{sec:tabular-fail}
% ===========================================================

\subsection{Проклятие размерности}

Чтобы понять необходимость аппроксимации функций, нужно сначала количественно оценить, как резко растёт пространство состояний с усложнением задачи.  В табличном МПР с $|\cS|$ состояниями и $|\cA|$ действиями представление функции ценности действия $Q : \cS \times \cA \to \R$ требует хранения $|\cS| \cdot |\cA|$ чисел с плавающей точкой.  Для небольшого клеточного мира со $100$ клетками и $4$ действиями это $400$ чисел\,---\,тривиально.  Но размер пространства состояний не является линейной функцией входной размерности задачи; он экспоненциален.

\begin{definition}[Проклятие размерности]
\label{def:curse-dim}
\index{curse of dimensionality}
Пусть состояние описывается $n$ признаками, каждый принимает $k$ дискретных значений.  Общее число состояний равно $k^n$.  Следовательно, затраты памяти и выборочная сложность табличных методов растут \emph{экспоненциально} по числу признаков $n$.  Это явление называется \textbf{проклятием размерности}.
\end{definition}

Три конкретных примера делают это наглядным.

\begin{example}[Видеоигры Atari]
\label{ex:atari-states}
\index{Atari}
Atari 2600 производит кадры $210 \times 160$ пикселей с $128$ цветами на пиксель.  Мних и др. (2015) предобрабатывают кадры до $84 \times 84$ градаций серого и укладывают четыре последовательных кадра для захвата движения, давая тензор наблюдения с $84 \times 84 \times 4 = 28\,224$ значениями, каждое --- целое число в $\{0, \ldots, 255\}$.  Число различных наблюдений составляет не более $256^{28\,224}$ --- число с примерно $68\,000$ десятичными цифрами.  Табличная $Q$-функция на этом пространстве не просто непрактична; она концептуально абсурдна.  Тем не менее DQN, используя свёрточную нейронную сеть с $\sim 1.5$ млн параметров, достиг сверхчеловеческой производительности во многих из этих игр.  Сеть преуспевает именно потому, что обобщает по визуально схожим кадрам.
\end{example}

\begin{example}[Непрерывная робототехника]
\label{ex:robotics-states}
\index{robotics!state space}
Симулированный гуманоидный робот (например, среда MuJoCo Humanoid-v4) имеет вектор состояния $s \in \R^{376}$, включающий угловые положения суставов, угловые скорости и контактные силы.  Это пространство состояний не просто велико\,---\,оно непрерывно и несчётно бесконечно.  Никакая схема дискретизации не может свести его к управляемой таблице без серьёзных ошибок приближения.  Более того, пространство действий --- $\R^{17}$, поэтому операция $\argmax$, центральная для Q-learning, становится задачей непрерывной оптимизации на каждом шаге.
\end{example}

\begin{example}[Язык и текст]
\label{ex:language-states}
\index{language models!state space}
В приложениях ОП к дообучению языковых моделей (подробнее в главе~\ref{ch:practical-rlhf}) состояние --- это текущий префикс токенов\,---\,вся история диалога до текущего шага.  При словарном запасе размером $50\,000$ и контекстном окне в $2\,048$ токенов число возможных префиксов астрономически велико.  Более того, семантика двух префиксов, отличающихся одним словом, может быть почти идентичной, поэтому любое полезное представление должно достигать семантического обобщения, а не только синтаксического.
\end{example}

\subsection{Аргумент обобщения}

Помимо чисто проблемы хранения, табличные методы имеют фундаментальную статистическую неэффективность: каждое состояние рассматривается как полностью отдельный объект.  Обучение тому, что ценность состояния $s$ высока, не даёт никакой информации о состоянии $s'$, даже если $s$ и $s'$ визуально идентичны за исключением положения одного пикселя.  Говоря языком статистики, табличные методы обладают нулевым \emph{обобщением}: каждое значение должно оцениваться независимо, по переходам, посещающим конкретное состояние.

В больших пространствах состояний большинство состояний посещается не более одного раза или вообще не посещается в ходе разумного обучающего прогона.  Табличные методы присваивают таким состояниям их начальные значения (часто ноль) на неопределённый срок.  Аппроксимация функций устраняет это, разделяя статистическую мощь между состояниями: посещение состояния $s$ неявно обновляет оценочные ценности всех состояний $s'$, для которых $\varphi(s')$ похоже на $\varphi(s)$, где $\varphi$ обозначает признаковое представление.  Эта \emph{неявная регуляризация} является основной пользой параметрических функций ценности.

% ===========================================================
\section{Линейная аппроксимация функций ценности}
\label{sec:linear-vfa}
% ===========================================================

\subsection{Признаковые векторы и линейный анзац}

Простейшее параметрическое приближение представляет функцию ценности как линейную комбинацию фиксированных признаков.

\begin{definition}[Линейная аппроксимация функций ценности]
\label{def:linear-vfa}
\index{linear value function approximation}
Пусть $\varphi : \cS \to \R^d$ --- \textbf{признаковое отображение}\index{feature map}, назначающее каждому состоянию признаковый вектор.  \textbf{Линейная аппроксимация функций ценности} имеет вид:
\begin{equation}
\label{eq:linear-vfa}
V_\mathbf{w}(s) = \mathbf{w}^\top \varphi(s) = \sum_{i=1}^d w_i \varphi_i(s),
\end{equation}
где $\mathbf{w} \in \R^d$ --- вектор весов.  Для функций ценности действия используются признаки $\varphi(s, a) \in \R^d$ и записывается $Q_\mathbf{w}(s, a) = \mathbf{w}^\top \varphi(s, a)$.
\end{definition}

Линейный анзац привлекателен по трём причинам.  Во-первых, это простейшая возможная параметризация: приближение однозначно определяется $d$ вещественными числами вне зависимости от $|\cS|$.  Во-вторых, многие интересные признаковые отображения могут быть спроектированы вручную\,---\,тайловые кодирования, радиальные базисные функции, полиномиальные базисы, фурье-признаки\,---\,каждое кодирует конкретные априорные убеждения о структуре функции ценности.  В-третьих, и самое главное, линейная аппроксимация функций является постановкой, в которой сходимость TD можно установить с математической строгостью, как будет показано в разделе~\ref{sec:projected-bellman}.

\begin{example}[Полиномиальные признаки для одномерного коридора]
\label{ex:poly-features}
\index{polynomial features}
Рассмотрим коридор из $N = 1\,000$ дискретных позиций.  В табличном случае $V^\pi$ требует $N = 1\,000$ параметров.  Если, однако, истинная функция ценности гладкая\,---\,что вероятно, если вознаграждение сосредоточено на одном конце и коэффициент дисконтирования умеренный\,---\,мы можем аппроксимировать её полиномиальным базисом:
\[
\varphi(s) = \bigl(1,\; s/N,\; (s/N)^2,\; (s/N)^3\bigr)^\top \in \R^4.
\]
Это сокращает число параметров с $1\,000$ до $4$, ценой невозможности точного представления быстро меняющихся функций ценности.  Ошибка приближения, возникающая при ограничении данным классом функций, называется \textbf{ошибкой представления}\index{representation error}, и она неизбежна всякий раз, когда истинная функция ценности лежит вне семейства приближений.
\end{example}

\subsection{Тайловое кодирование}

Особенно полезным признаковым отображением для непрерывных пространств состояний является \emph{тайловое кодирование}%
\index{tile coding}, разработанное Саттоном (1996).  Пространство состояний $\R^n$ покрывается несколькими перекрывающимися сетками, называемыми \emph{тайлингами}.  Для данного состояния $s$ ровно один тайл в каждом тайлинге активен, давая двоичный признаковый вектор, ненулевые элементы которого указывают, какой тайл активен в каждом тайлинге.  Если имеется $m$ тайлингов, каждый разбивающий пространство на $k$ тайлов, признаковый вектор имеет размерность $mk$ и содержит ровно $m$ единиц.  Результирующая линейная аппроксимация эквивалентна таблице значений тайлов, но перекрывающаяся структура тайлингов обеспечивает гладкую интерполяцию между близкими состояниями.

\subsection{Линейное TD(0)}
\label{subsec:linear-td}

Вспомним из главы~\ref{ch:td}, что TD(0) для оценивания стратегии поддерживает оценку ценности и обновляет её после каждого перехода $(S_t, A_t, R_{t+1}, S_{t+1})$, используя ошибку TD:
\[
\delta_t = R_{t+1} + \gamma V(S_{t+1}) - V(S_t).
\]
Когда $V$ параметризована как $V_\mathbf{w}(s) = \mathbf{w}^\top \varphi(s)$, естественным расширением является обновление вектора весов в направлении, уменьшающем квадрат ошибки TD.  Это даёт обновление \textbf{полуградиентного TD(0)}\index{semi-gradient TD(0)} (подробно анализируется в разделе~\ref{sec:semi-gradient}):
\begin{equation}
\label{eq:linear-td-update}
\mathbf{w} \leftarrow \mathbf{w} + \alpha\, \delta_t\, \varphi(S_t),
\end{equation}
где $\delta_t = R_{t+1} + \gamma \mathbf{w}^\top \varphi(S_{t+1}) - \mathbf{w}^\top \varphi(S_t)$.
Заметим, что в линейном случае $\nabla_\mathbf{w} V_\mathbf{w}(S_t) = \varphi(S_t)$, поэтому~\eqref{eq:linear-td-update} является шагом полуградиента, применённым к $V_\mathbf{w}$.

Критическое свойство этого обновления состоит в том, что изменение $\mathbf{w}$ затрагивает \emph{все состояния одновременно}: предсказание ценности в состоянии $s'$ изменяется на $\alpha \delta_t \, \varphi(s') \cdot \varphi(S_t)$, пропорционально схожести $s'$ с $S_t$ по признаковому скалярному произведению.  Это и есть механизм обобщения в линейной аппроксимации функций.

\begin{theorem}[Сходимость линейного TD(0) при выборке на текущей стратегии]
\label{thm:linear-td-convergence}
\index{linear TD(0)!convergence}
Пусть $\pi$ --- фиксированная стратегия и $\mu$ --- её стационарное распределение по $\cS$ (предполагается единственным и положительным везде).  Предположим, что признаковая матрица $\Phi \in \R^{|\cS| \times d}$ имеет полный ранг по столбцам.  Если выборки $(S_t, R_{t+1}, S_{t+1})$ берутся на текущей стратегии (т.е. $S_t \sim \mu$ в установившемся режиме) и шаги удовлетворяют $\sum_t \alpha_t = \infty$ и $\sum_t \alpha_t^2 < \infty$, то итерации $\mathbf{w}_t$ линейного TD(0) сходятся почти наверное к единственному вектору весов $\mathbf{w}_{\mathrm{TD}}$, удовлетворяющему проецированному уравнению Беллмана $\Phi \mathbf{w}_{\mathrm{TD}} = \Pi_\mu T^\pi (\Phi \mathbf{w}_{\mathrm{TD}})$ (определяется в разделе~\ref{sec:projected-bellman}).
\end{theorem}

Доказательство теоремы~\ref{thm:linear-td-convergence} опирается на теорию стохастической аппроксимации и свойства сжатия проецированного оператора Беллмана, которые мы сейчас развиваем.

% ===========================================================
\section{Проецированное уравнение Беллмана}
\label{sec:projected-bellman}
% ===========================================================

\subsection{Геометрия аппроксимации функций}

Чтобы понять, к чему в действительности сходится линейное TD(0), нужно думать геометрически.  Линейный класс функций $\mathcal{F} = \{\Phi \mathbf{w} : \mathbf{w} \in \R^d\}$ --- это $d$-мерное подпространство $\R^{|\cS|}$ (отождествляя функции ценности с векторами, проиндексированными состояниями).  Оператор Беллмана $T^\pi$ отображает любую функцию ценности $v \in \R^{|\cS|}$ в новую функцию $T^\pi v$, но в общем случае $T^\pi v \notin \mathcal{F}$ даже при $v \in \mathcal{F}$.  Следовательно, TD-методы, работающие в $\mathcal{F}$, не могут следовать $T^\pi$ напрямую.

Решение --- \emph{проецировать} $T^\pi v$ обратно в $\mathcal{F}$ после каждого шага Беллмана.

\begin{definition}[$\mu$-взвешенное скалярное произведение и проекция]
\label{def:mu-projection}
\index{projection operator!$\mu$-weighted}
Пусть $\mu : \cS \to (0, \infty)$ --- положительное взвешивание (как правило, стационарное распределение $\pi$ на текущей стратегии).  Определим \textbf{$\mu$-взвешенную норму}\index{$\mu$-norm}:
\begin{equation}
\label{eq:mu-norm}
\|v\|_\mu^2 = \sum_{s \in \cS} \mu(s)\, v(s)^2 = \mathbf{v}^\top D_\mu \mathbf{v},
\end{equation}
где $D_\mu = \mathrm{diag}(\mu)$.  \textbf{Ортогональная проекция} $v$ на $\mathcal{F}$ в этой норме:
\begin{equation}
\label{eq:projection-def}
\Pi_\mu v = \Phi\bigl(\Phi^\top D_\mu \Phi\bigr)^{-1} \Phi^\top D_\mu v
= \argmin_{f \in \mathcal{F}} \|f - v\|_\mu.
\end{equation}
\end{definition}

Матрица $\Pi_\mu = \Phi(\Phi^\top D_\mu \Phi)^{-1}\Phi^\top D_\mu$ является ортогональным проектором в $\mu$-взвешенном скалярном произведении.  Она идемпотентна ($\Pi_\mu^2 = \Pi_\mu$) и удовлетворяет $\|I - \Pi_\mu\|_\mu \leq 1$, то есть проекция может только приблизить вектор \emph{к} $\mathcal{F}$, но не удалить его.

\begin{definition}[Проецированное уравнение Беллмана]
\label{def:projected-bellman}
\index{projected Bellman equation}
\textbf{Проецированное уравнение Беллмана} для стратегии $\pi$:
\begin{equation}
\label{eq:projected-bellman}
\hat{V} = \Pi_\mu T^\pi \hat{V}, \qquad \hat{V} \in \mathcal{F}.
\end{equation}
Решение $\hat{V} = \Phi \mathbf{w}_{\mathrm{TD}}$ называется \textbf{неподвижной точкой TD}\index{TD fixed point}.
\end{definition}

Проецированное уравнение Беллмана имеет ясную геометрическую интерпретацию: найти функцию в классе приближений $\mathcal{F}$, которая после одного шага Беллмана и повторного проецирования возвращает к себе.  Это наилучшее, что мы можем ожидать, когда точная согласованность Беллмана недостижима в $\mathcal{F}$.

\begin{example}[Проецированное уравнение Беллмана для трёх состояний]
\label{ex:three-state-projection}
Предположим $|\cS| = 3$ с равномерным взвешиванием $\mu = (1/3, 1/3, 1/3)$, и признаковая матрица:
\[
\Phi = \begin{pmatrix} 1 & 0 \\ 0 & 1 \\ 1 & 1 \end{pmatrix} \in \R^{3 \times 2}.
\]
Семейство приближений $\mathcal{F}$ имеет размерность $2$, поэтому значения во всех трёх состояниях не могут задаваться независимо: $V(s_3) = V(s_1) + V(s_2)$ всегда.  Проекция $\Pi_\mu v$ любого $v \in \R^3$ находит веса $\mathbf{w}^* = \argmin_\mathbf{w} \|v - \Phi \mathbf{w}\|_\mu^2$, решая $\Phi^\top D_\mu \Phi \mathbf{w}^* = \Phi^\top D_\mu v$, что даёт:
\[
\mathbf{w}^* = \tfrac{1}{2}\begin{pmatrix} v_1 + v_3 - v_2 \\ v_2 + v_3 - v_1 \end{pmatrix},
\]
то есть проекция «делит вклад $s_3$» между двумя базисными направлениями.
\end{example}

\subsection{Существование неподвижной точки TD}

\begin{theorem}[Существование и единственность неподвижной точки TD]
\label{thm:td-fixed-point}
\index{TD fixed point!existence}
При допущениях теоремы~\ref{thm:linear-td-convergence} проецированный оператор Беллмана $\Pi_\mu T^\pi : \mathcal{F} \to \mathcal{F}$ является сжатием в $\mu$-норме с коэффициентом $\gamma$:
\begin{equation}
\label{eq:projected-contraction}
\|\Pi_\mu T^\pi u - \Pi_\mu T^\pi v\|_\mu \leq \gamma \|u - v\|_\mu
\qquad \forall\, u, v \in \mathcal{F}.
\end{equation}
Следовательно, проецированное уравнение Беллмана~\eqref{eq:projected-bellman} имеет единственное решение $\hat{V} = \Phi \mathbf{w}_{\mathrm{TD}} \in \mathcal{F}$.
\end{theorem}

\begin{proof}
Сам оператор Беллмана $T^\pi$ является сжатием в $\mu$-норме с коэффициентом $\gamma$ (это следует из свойства сжатия $T^\pi$ в супремум-норме и эквивалентности этих норм при положительном $\mu$).  Формально, для любых $u, v$:
\[
\|T^\pi u - T^\pi v\|_\mu \leq \gamma \|u - v\|_\mu.
\]
Поскольку $\Pi_\mu$ --- неэкспансивная проекция ($\|\Pi_\mu x\|_\mu \leq \|x\|_\mu$ для всех $x$), имеем:
\[
\|\Pi_\mu T^\pi u - \Pi_\mu T^\pi v\|_\mu
\leq \|T^\pi u - T^\pi v\|_\mu
\leq \gamma \|u - v\|_\mu.
\]
Поскольку $\gamma < 1$ и $\mathcal{F}$ является полным метрическим пространством (замкнутым подпространством $\R^{|\cS|}$), теорема Банаха о неподвижной точке гарантирует единственную неподвижную точку.
\end{proof}

\subsection{Ошибка приближения в неподвижной точке TD}

Неподвижная точка TD $\hat{V}$ минимизирует определённую меру ошибки, но не минимизирует напрямую среднеквадратичную ошибку $\|\hat{V} - V^\pi\|_\mu^2$.  Следующая классическая оценка (Цицикис и Ван Рой, 1997) количественно определяет, насколько неподвижная точка TD отклоняется от наилучшего возможного приближения.

\begin{theorem}[Оценка ошибки приближения в неподвижной точке TD]
\label{thm:td-bound}
\index{TD fixed point!approximation error}
Пусть $\hat{V}^* = \Pi_\mu V^\pi$ --- наилучшее приближение $V^\pi$ в $\mathcal{F}$.  Тогда:
\begin{equation}
\label{eq:td-bound}
\|\hat{V} - V^\pi\|_\mu \leq \frac{1}{\sqrt{1 - \gamma^2}}\, \|\hat{V}^* - V^\pi\|_\mu.
\end{equation}
\end{theorem}

Коэффициент $1/\sqrt{1-\gamma^2}$ --- это «цена» бутстрэппинга: неподвижная точка TD может быть в $1/\sqrt{1-\gamma^2}$ раз хуже наилучшего приближения в классе.  При $\gamma = 0.99$ этот коэффициент составляет примерно $7.1$, иллюстрируя, что бутстрэппинг может существенно усилить ошибки приближения в задачах с длинными горизонтами времени.

% ===========================================================
\section{Методы полуградиента}
\label{sec:semi-gradient}
% ===========================================================

\subsection{Истинный градиентный спуск для предсказания}

Прежде чем изучать методы полуградиента, поучительно понять, как выглядел бы истинный градиентный спуск для задачи предсказания.  Определим \textbf{среднеквадратичную ошибку ценности}\index{mean-squared value error}:
\begin{equation}
\label{eq:msve}
\overline{\mathrm{VE}}(\mathbf{w})
= \|V^\pi - V_{\mathbf{w}}\|_{\mu}^{2}
= \sum_{s \in \cS} \mu(s)\bigl[V^\pi(s) - V_\mathbf{w}(s)\bigr]^2,
\end{equation}
где $\mu$ --- распределение на текущей стратегии.  Истинный градиент $\overline{\mathrm{VE}}$ по $\mathbf{w}$:
\begin{equation}
\label{eq:true-grad-msve}
-\frac{1}{2}\nabla_\mathbf{w}\overline{\mathrm{VE}}(\mathbf{w})
= \sum_s \mu(s)\bigl[V^\pi(s) - V_\mathbf{w}(s)\bigr]\nabla_\mathbf{w} V_\mathbf{w}(s).
\end{equation}
К сожалению, $V^\pi(s)$ неизвестна, поэтому мы не можем вычислить этот градиент напрямую.  Методы Монте-Карло подставляют выборочную отдачу $G_t \approx V^\pi(S_t)$, давая обновление:
\[
\mathbf{w} \leftarrow \mathbf{w} + \alpha\bigl[G_t - V_\mathbf{w}(S_t)\bigr]\nabla_\mathbf{w} V_\mathbf{w}(S_t).
\]
Это истинный стохастический шаг градиентного спуска, потому что $\E[G_t \mid S_t] = V^\pi(S_t)$ и градиент $\nabla_\mathbf{w} V_\mathbf{w}(S_t)$ не входит в цель $G_t$.

\subsection{Почему TD нарушает истинный градиентный спуск}

Методы TD заменяют $G_t$ одношаговой целью бутстрэппинга $R_{t+1} + \gamma V_\mathbf{w}(S_{t+1})$.  Полученное обновление:
\begin{equation}
\label{eq:semi-grad-td}
\mathbf{w} \leftarrow \mathbf{w} + \alpha\underbrace{\bigl[R_{t+1} + \gamma V_\mathbf{w}(S_{t+1}) - V_\mathbf{w}(S_t)\bigr]}_{\delta_t}\nabla_\mathbf{w} V_\mathbf{w}(S_t).
\end{equation}
Это \emph{похоже} на шаг градиентного спуска, но таковым не является.

\begin{definition}[Метод полуградиента]
\label{def:semi-gradient}
\index{semi-gradient method}
\textbf{Метод полуградиента} --- это метод, в котором направление обновления вычисляется так, как если бы $\delta_t$ был сигналом ошибки, не зависящим от $\mathbf{w}$, то есть градиент берётся только через $V_\mathbf{w}(S_t)$, но не через цель $V_\mathbf{w}(S_{t+1})$.  Обновление~\eqref{eq:semi-grad-td} является шагом полуградиента, поскольку игнорируется градиент $\delta_t$ по $\mathbf{w}$ через его зависимость от $V_\mathbf{w}(S_{t+1})$.
\end{definition}

Чтобы понять, почему градиент через $V_\mathbf{w}(S_{t+1})$ опускается, заметим, что если бы мы вычисляли истинный градиент $\frac{1}{2}\delta_t^2$, то получили бы:
\[
\nabla_\mathbf{w}\tfrac{1}{2}\delta_t^2
= \delta_t\bigl(\gamma\nabla_\mathbf{w}V_\mathbf{w}(S_{t+1}) - \nabla_\mathbf{w}V_\mathbf{w}(S_t)\bigr).
\]
Полуградиент TD отбрасывает слагаемое $\gamma\nabla_\mathbf{w}V_\mathbf{w}(S_{t+1})$.  Это означает, что направление обновления \emph{не является} отрицательным градиентом никакой фиксированной целевой функции.  Как следствие, стандартные теоремы о сходимости стохастического градиентного спуска не применяются непосредственно.

\begin{remark}
\label{rem:semi-grad-motivation}
Опущение градиента бутстрэппинга является намеренным, а не ошибкой.  Если бы включить его, то полученный алгоритм \emph{остаточного градиента} сходился бы в среднеквадратичной ошибке Беллмана, но страдал бы от проблемы «двойной выборки»: для вычисления полного градиента требуются две независимые выборки $S_{t+1}$ для одной и той же пары $(S_t, A_t)$, что обычно недоступно в едином потоке взаимодействия.  Методы полуградиента избегают этого, рассматривая цель как фиксированную; их практические свойства сходимости, как правило, лучше, несмотря на отсутствие гарантии градиентного спуска.
\end{remark}

\subsection{Алгоритм полуградиентного TD}

\begin{algorithm}[ht]
\caption{Полуградиентное TD(0) для оценивания стратегии}
\label{alg:semi-grad-td0}
\KwIn{Стратегия $\pi$, признаковое отображение $\varphi$, шаг $\alpha > 0$, дисконт $\gamma \in [0,1)$}
Инициализировать $\mathbf{w} \in \R^d$ (например, $\mathbf{w} \leftarrow \mathbf{0}$)\;
\For{каждый эпизод}{
  Инициализировать состояние $S_0$\;
  \For{$t = 0, 1, 2, \ldots$ пока $S_t$ не терминальное}{
    Выбрать $A_t \sim \pi(\cdot \mid S_t)$\;
    Наблюдать $R_{t+1}$ и $S_{t+1}$\;
    $\delta_t \leftarrow R_{t+1} + \gamma \mathbf{w}^\top \varphi(S_{t+1}) - \mathbf{w}^\top \varphi(S_t)$\;
    $\mathbf{w} \leftarrow \mathbf{w} + \alpha\, \delta_t\, \varphi(S_t)$\;
  }
}
\end{algorithm}

\subsection{Сходимость полуградиентного TD}

Для линейного случая полуградиентное TD(0) при выборке на текущей стратегии сходится к неподвижной точке TD $\mathbf{w}_{\mathrm{TD}}$, характеризуемой теоремой~\ref{thm:td-fixed-point}.  Аргумент опирается на \emph{метод обыкновенного дифференциального уравнения (ОДУ)}\index{ODE method}: показывается, что поле среднего обновления
\[
\bar{f}(\mathbf{w}) = \E_\mu[\delta_t \varphi(S_t) \mid \mathbf{w}]
= \Phi^\top D_\mu (T^\pi(\Phi\mathbf{w}) - \Phi\mathbf{w})
\]
является сжатием к $\mathbf{w}_{\mathrm{TD}}$, а затем применяется стандартная теорема (Кушнер и Йинь, 2003), утверждающая, что стохастическая аппроксимация сходится к равновесию соответствующего ОДУ.

Для \emph{нелинейных} аппроксиматоров функций даже полуградиентное TD на текущей стратегии не гарантировано сходится.  На практике TD-методы на основе нейронных сетей часто сходятся к полезным решениям, однако контрпримеры существуют.  Сходимость в нелинейном случае остаётся активной областью исследований.

% ===========================================================
\section{Смертельная триада}
\label{sec:deadly-triad}
% ===========================================================

\subsection{Три ингредиента нестабильности}

Гарантия сходимости теоремы~\ref{thm:linear-td-convergence} критически зависит от слов «выборка на текущей стратегии».  При нарушении любого из трёх условий хорошо поведённый ландшафт линейного TD разрушается.  При одновременном нарушении всех трёх даже линейная аппроксимация функций может расходиться.  Саттон и Барто (2018) называют эту опасную комбинацию \textbf{смертельной триадой}\index{deadly triad}.

\begin{definition}[Смертельная триада]
\label{def:deadly-triad}
\index{deadly triad}
\textbf{Смертельная триада} означает одновременное присутствие трёх элементов:
\begin{enumerate}
\item \textbf{Аппроксимация функций:} использование параметрической функции $V_\mathbf{w}$ (или $Q_\mathbf{w}$) с $d \ll |\cS|$ вместо отдельной табличной записи на состояние.
\item \textbf{Бутстрэппинг:} формирование целей из текущей оценки, как в TD-методах, вместо ожидания полной отдачи, как в Монте-Карло.
\item \textbf{Обучение на другой стратегии:} обновление параметров по данным, сгенерированным \emph{поведенческой стратегией} $b$, отличной от оцениваемой или улучшаемой \emph{целевой стратегии} $\pi$.
\end{enumerate}
Сочетание всех трёх может вызвать расходимость параметров даже при линейных аппроксиматорах и даже при бесконечном объёме данных.
\end{definition}

\begin{figure}[ht]
\centering
\begin{tikzpicture}[
  every node/.style={font=\small},
  tri/.style={draw, fill=gray!15, regular polygon, regular polygon sides=3,
              minimum size=5.5cm, inner sep=0pt},
  label/.style={font=\small\bfseries, align=center}
]
  % Треугольник смертельной триады
  \node[tri] (T) {};
  % Вершины
  \coordinate (A) at (T.corner 1);  % сверху
  \coordinate (B) at (T.corner 2);  % снизу слева
  \coordinate (C) at (T.corner 3);  % снизу справа

  \node[label, above=0.45cm of A, align=center] {Аппроксимация\\функций};
  \node[label, below=0.25cm of B, align=right] {Бутстрэппинг};
  \node[label, below=0.25cm of C, align=left] {Обучение на другой\\стратегии};

  \node[font=\itshape\small, align=center] at ($(A)!0.5!(B)!0.33!(C)$)
    {Потенциальная\\расходимость};

  % Аннотации безопасных зон на рёбрах (сдвинуты внутрь треугольника)
  \node[font=\tiny, sloped, above, inner sep=1pt] at ($(A)!0.5!(B)$) {МК (без бутстрэппинга)};
  \node[font=\tiny, above, inner sep=2pt] at ($(B)!0.5!(C)+(0,0.1)$) {Таблица (без аппрокс.)};
  \node[font=\tiny, sloped, above, inner sep=1pt] at ($(A)!0.5!(C)$) {TD на текущей стратегии};
\end{tikzpicture}
\caption{Смертельная триада.  Каждое ребро треугольника «безопасно»: два из трёх ингредиентов вместе управляемы.  Только когда все три вершины активны одновременно, нестабильность становится неизбежной в общем случае.  Подписи рёбер указывают, какой элемент отсутствует вдоль каждой стороны.}
\label{fig:deadly-triad}
\end{figure}

Важно понимать, что каждая пара ингредиентов управляема:
\begin{itemize}[leftmargin=*]
\item \emph{Аппроксимация функций + бутстрэппинг, на текущей стратегии:} Линейное TD(0) сходится к неподвижной точке TD (теорема~\ref{thm:linear-td-convergence}).
\item \emph{Аппроксимация функций + другая стратегия, без бутстрэппинга:} Монте-Карло на другой стратегии с линейной аппроксимацией функций сходится, поскольку градиент является истинным градиентом среднеквадратичной ошибки ценности.
\item \emph{Бутстрэппинг + другая стратегия, табличное:} Q-learning сходится к $Q^*$ в табличном случае (глава~\ref{ch:td}, теорема о сходимости Q-learning).
\end{itemize}

Проблема возникает только при сочетании всех трёх.

\subsection{Контрпример Бэрда}
\label{subsec:baird}

Наиболее ясная иллюстрация смертельной триады --- \emph{контрпример Бэрда}\index{Baird's counterexample} (Бэрд, 1995).  Рассмотрим простой МПР с $7$ состояниями и $2$ действиями.  Шесть состояний расположены по схеме звезды; седьмое --- центральный узел.  Поведенческая стратегия всегда выбирает первое действие (переводящее в одно из шести состояний-лучей, равномерно) с вероятностью $\frac{6}{7}$ и второе действие (переводящее в узел) с вероятностью $\frac{1}{7}$.  Целевая стратегия всегда выбирает второе действие.  Все вознаграждения равны нулю, $\gamma = 0.99$.

Признаковые векторы выбраны так, что семейство приближений линейно.  Признаковый вектор для состояния-узла $s_7$ --- $(0,0,0,0,0,0,1,2)^\top$, а для состояния-луча $s_i$ ($i = 1, \ldots, 6$) --- вектор с $2$ на позиции $i$ и $1$ на позиции $8$.  Вектор весов $\mathbf{w} \in \R^8$ инициализируется значением $(1,1,1,1,1,1,10,1)^\top$.

\begin{theorem}[Контрпример Бэрда]
\label{thm:baird}
\index{Baird's counterexample!divergence}
При полуградиентном TD(0) на другой стратегии с описанными выше признаковыми векторами вектор весов $\mathbf{w}_t$ расходится к бесконечности почти наверное, несмотря на то, что признаковая матрица имеет полный ранг по столбцам и шаг удовлетворяет стандартным условиям Роббинса--Монро.
\end{theorem}

Механизм расходимости тонок: распределение на другой стратегии $\mu$ концентрирует слишком много веса на состояниях-лучах (которые часто посещаются при поведенческой стратегии), тогда как обновления Беллмана преимущественно «толкают» ценности этих состояний вверх, поскольку целевая стратегия предпочитает идти в узел.  Проекция обратно в класс функций затем усиливает этот толчок вместо его демпфирования, создавая положительную обратную связь.

\begin{remark}[Что нас учит контрпример Бэрда]
\label{rem:baird-lesson}
Контрпример Бэрда --- не патологическое курьёзное явление; это фундаментальный результат об ограничениях TD-методов.  Он говорит нам, что сочетание полуградиентного бутстрэппинга и данных с другой стратегии изначально опасно, независимо от того, насколько хорошо ведёт себя признаковое представление.  Решения, исследуемые на практике\,---\,коррекции взвешиванием по важности, методы градиентного TD (TDC, GTD2) и инженерные стабилизации DQN\,---\,представляют различные стратегии выхода из опасностей триады без потери преимуществ данных с другой стратегии.
\end{remark}

% ===========================================================
\section{Deep Q-Networks}
\label{sec:dqn}
% ===========================================================

\subsection{От линейного к нелинейному приближению}

Успехи линейной аппроксимации функций реальны, но ограничены: метод требует вручную спроектированных признаков, и в высокомерных перцептивных областях, таких как сырые пиксели изображений, ни один человеческий разработчик не может заранее определить правильные признаки.  Следующий естественный шаг --- использовать глубокую нейронную сеть\index{deep neural network} в качестве аппроксиматора, позволяя сети обучать своё собственное внутреннее представление из данных.

Проблема в том, что глубокие сети выявляют все трудности смертельной триады в полной мере.  При Q-learning с нейронной сетью $Q_\theta$ и жадным выбором действий цель для пары состояние--действие $(S_t, A_t)$:
\begin{equation}
\label{eq:naive-deep-q-target}
y_t = R_{t+1} + \gamma \max_{a'} Q_\theta(S_{t+1}, a'),
\end{equation}
а минимизируемая функция потерь --- $(y_t - Q_\theta(S_t, A_t))^2$.  Но и предсказание $Q_\theta(S_t, A_t)$, и цель $y_t$ зависят от одних и тех же параметров $\theta$, поэтому каждый шаг градиента одновременно изменяет предсказание \emph{и} цель.  Это создаёт проблему движущейся цели: сеть гонится за величиной, которая меняется с каждым обновлением, приводя на практике к осцилляциям и расходимости.

Две дополнительные патологии усугубляют проблему.  Переходы $(S_t, A_t, R_{t+1}, S_{t+1})$, собираемые последовательным агентом, \emph{сильно коррелированы}: последовательные переходы разделяют одну и ту же последовательность состояний.  Обучение нейронной сети на коррелированных минипакетах нарушает предположение о н.о.р. стохастического градиентного спуска и заставляет сеть переобучаться на недавнем опыте, «забывая» более раннее обучение.  Кроме того, распределение обучающих данных меняется по мере улучшения стратегии, делая ландшафт оптимизации нестационарным.

\subsection{Буфер повторного воспроизведения опыта}
\label{subsec:experience-replay}

Первое стабилизирующее нововведение, введённое Мних и др. (2013, 2015), --- \textbf{повторное воспроизведение опыта}\index{experience replay}.  По мере взаимодействия агента со средой каждый переход $(S_t, A_t, R_{t+1}, S_{t+1})$ сохраняется в \emph{буфере повторного воспроизведения}\index{replay buffer} $\mathcal{D}$ ёмкостью $N$.  Во время обучения минипакеты переходов извлекаются \emph{равномерно случайно} из $\mathcal{D}$:
\begin{equation}
\label{eq:replay-sampling}
\{(s_i, a_i, r_i, s'_i)\}_{i=1}^B \sim \mbox{Равном.}(\mathcal{D}).
\end{equation}
Это простое изменение даёт два важных эффекта.  Во-первых, оно \emph{декоррелирует} обучающие данные: последовательные минипакеты извлекаются из многих различных временных шагов, разрывая временную корреляцию, дестабилизирующую онлайн-обновления.  Во-вторых, каждый переход используется потенциально для многих шагов градиентного спуска, повышая эффективность использования данных.  Цена --- необходимость хранить большой буфер в памяти и введение неявного компонента другой стратегии: переходы, сохранённые давно, были сгенерированы более ранней, другой стратегией.

\subsection{Целевые сети}
\label{subsec:target-networks}

Второе стабилизирующее нововведение --- \textbf{целевая сеть}\index{target network}.  DQN поддерживает две сети с одинаковой архитектурой: \emph{онлайн-сеть} $Q_\theta$ (также называемая \emph{поведенческой сетью}) и \emph{целевая сеть} $Q_{\theta^-}$.  Онлайн-сеть обновляется на каждом шаге градиентного спуска.  Параметры целевой сети $\theta^-$ фиксируются на $C$ шагов (обычно $C = 10\,000$), а затем синхронизируются с онлайн-сетью: $\theta^- \leftarrow \theta$.

С целевой сетью функция потерь принимает вид:
\begin{equation}
\label{eq:dqn-loss}
L(\theta) = \E_{(s,a,r,s') \sim \mathcal{D}}\!\left[
\Bigl(r + \gamma \max_{a'} Q_{\theta^-}(s', a') - Q_\theta(s, a)\Bigr)^2
\right].
\end{equation}
Градиент берётся только по $\theta$, при $\theta^-$ как константе.  Фиксируя цель на $C$ шагов, метод снижает корреляцию между обновляемыми параметрами и преследуемой ими целью.  Цель не является полностью стационарной\,---\,она всё равно меняется каждые $C$ шагов\,---\,но на практике это существенно стабилизирует обучение.


\begin{figure}[ht]
\centering
\begin{tikzpicture}[scale=0.9, transform shape,
  box/.style={draw, rounded corners, minimum width=2.8cm, minimum height=0.9cm,
              font=\small, fill=gray!10, align=center},
  buf/.style={draw, rounded corners, minimum width=2.8cm, minimum height=0.9cm,
              font=\small, fill=blue!10, align=center},
  lossbox/.style={draw, rounded corners, minimum width=2.8cm, minimum height=0.9cm,
              font=\small, fill=orange!10, align=center},
  arr/.style={-{Latex[length=2mm]}, thick},
  darr/.style={-{Latex[length=2mm]}, thick, dashed}
]

  % Строка 1: Онлайн-сеть (слева) и Целевая сеть (справа)
  \node[box] (online) at (0,0) {Онлайн-сеть\\$Q_\theta$};
  \node[box] (target) at (5.5,0) {Целевая сеть\\$Q_{\theta^-}$};

  % Строка 2: Функция потерь по центру
  \node[lossbox] (loss) at (2.75,-2.8) {Потери $L(\theta)$};

  % Строка 3: Буфер повторного воспроизведения
  \node[buf] (buf) at (2.75,-5.0) {Буфер $\mathcal{D}$};

  % Пунктирная стрелка: копирование весов
  \draw[darr] (online.north) to[bend left=20]
    node[above, font=\scriptsize]{копировать каждые $C$ шагов} (target.north);

  % Онлайн-сеть -> Потери
  \draw[arr] (online.south) -- (loss.north west)
    node[midway, left, font=\scriptsize, xshift=-2pt]{предсказать $Q_\theta(s,a)$};

  % Целевая сеть -> Потери
  \draw[arr] (target.south) -- (loss.north east)
    node[midway, right, font=\scriptsize, xshift=2pt]{цель $\max_{a'}Q_{\theta^-\!}(s',a')$};

  % Буфер -> Потери
  \draw[arr] (buf.north) -- (loss.south)
    node[midway, right, font=\scriptsize, xshift=2pt]{минипакет};

  % Потери -> Онлайн-сеть (градиент)
  \draw[arr] (loss.west) to[out=180, in=250, looseness=1.2]
    node[left, font=\scriptsize, xshift=-4pt]{$\nabla_\theta L$}
    (online.south west);

\end{tikzpicture}
\caption{Архитектура DQN.  Онлайн-сеть $Q_\theta$ выбирает действия и обновляется на каждом шаге.  Целевая сеть $Q_{\theta^-}$ вычисляет TD-цель и синхронизируется каждые $C$ шагов.  Переходы хранятся в буфере повторного воспроизведения $\mathcal{D}$ и извлекаются случайными минипакетами для обучения.  Пунктирная стрелка: периодическое жёсткое копирование $\theta^- \leftarrow \theta$.}
\label{fig:dqn-arch}
\end{figure}

\subsection{Алгоритм DQN}

\begin{algorithm}[ht]
\caption{Deep Q-Network (DQN)}
\label{alg:dqn}
\KwIn{Среда, сеть $Q_\theta$, ёмкость буфера $N$, размер пакета $B$, частота обновления $C$, расписание $\varepsilon$, $\gamma$, $\alpha$}
Инициализировать $\theta$ случайно; $\theta^- \leftarrow \theta$; $\mathcal{D} \leftarrow \emptyset$\;
\For{каждый эпизод}{
  Сбросить среду; наблюдать $s_0$\;
  \For{$t = 0, 1, 2, \ldots$}{
    С вероятностью $\varepsilon$: $a_t \leftarrow$ случайное действие; иначе $a_t \leftarrow \argmax_a Q_\theta(s_t, a)$\;
    Выполнить $a_t$; наблюдать $r_{t+1}$, $s_{t+1}$, флаг завершения\;
    Сохранить $(s_t, a_t, r_{t+1}, s_{t+1})$ в $\mathcal{D}$ (отбросить старейшее при заполнении)\;
    \If{$|\mathcal{D}| \geq B$}{
      Извлечь минипакет $\{(s_i, a_i, r_i, s'_i)\}_{i=1}^B \sim \mathcal{D}$\;
      \For{каждого перехода $i$}{
        \lIf{$s'_i$ терминальное}{$y_i \leftarrow r_i$}
        \lElse{$y_i \leftarrow r_i + \gamma \max_{a'} Q_{\theta^-}(s'_i, a')$}
      }
      Шаг градиента: $\theta \leftarrow \theta - \alpha \nabla_\theta \frac{1}{B}\sum_i (y_i - Q_\theta(s_i, a_i))^2$\;
      \If{$t \bmod C = 0$}{$\theta^- \leftarrow \theta$}
    }
    \lIf{завершено}{\textbf{break}}
  }
}
\end{algorithm}

\subsection{Архитектура и результаты на Atari}

Для игр Atari сеть DQN является свёрточной архитектурой: три свёрточных слоя (с $32$, $64$ и $64$ фильтрами размеров $8\times 8$, $4\times 4$, $3\times 3$ и шагами $4$, $2$, $1$ соответственно), за которыми следует полносвязный слой из $512$ нейронов и линейный выходной слой с одним нейроном на действие.  На вход подаётся стопка из четырёх $84\times 84$ кадров в градациях серого, предобработанных из сырого вывода $210\times 160$ RGB.

Мних и др. (2015) оценили DQN на $49$ играх Atari, используя \emph{единственную фиксированную архитектуру и гиперпараметры} для всех игр.  DQN достиг производительности выше уровня человека в $29$ из $49$ игр и установил новое состояние искусства почти во всех из них.  Принципиально важно, что DQN получает на вход только сырые пиксели и счёт игры, без какой-либо инженерии признаков, специфичной для игры.  Это впервые продемонстрировало, что единственный алгоритм глубокого ОП может достичь широкой компетентности в разнообразных задачах исключительно по высокоразмерному сенсорному вводу.

% ===========================================================
\section{Расширения DQN}
\label{sec:dqn-extensions}
% ===========================================================

Статья о DQN дала толчок волне алгоритмических улучшений.  Рассмотрим три, ставших стандартными компонентами современных систем глубокого ОП.

\subsection{Double DQN}
\label{subsec:double-dqn}

Тонкая, но важная патология Q-learning --- \textbf{смещение максимизации}\index{maximisation bias}: оператор $\max_{a'} Q(s', a')$ является завышенной оценкой $\max_{a'} Q^*(s', a')$, когда $Q$ содержит шум, поскольку максимум зашумлённых оценок смещён вверх.  Ван Хасселт и др. (2016) показали, что это смещение присутствует и вредоносно в DQN, заставляя агента переоценивать ценности и совершать неоптимальные действия.

Решение --- \textbf{Double DQN}\index{Double DQN} (DDQN): разделить выбор действия и оценку ценности, используя онлайн-сеть для выбора жадного действия и целевую сеть для его оценки:
\begin{equation}
\label{eq:ddqn-target}
y_t^{\mathrm{DDQN}} = R_{t+1} + \gamma Q_{\theta^-}\!\left(S_{t+1},\, \argmax_{a'} Q_\theta(S_{t+1}, a')\right).
\end{equation}
По сравнению с целью DQN в~\eqref{eq:dqn-loss} единственное изменение --- $\argmax$ и оценка $Q$ выполняются разными сетями.  Это простое изменение существенно снижает завышение оценок и улучшает производительность во многих играх Atari при пренебрежимых вычислительных накладных расходах.

\subsection{Дуэльный DQN}
\label{subsec:dueling-dqn}

Ван и др. (2016) заметили, что во многих состояниях ценность состояния $V^\pi(s)$ более информативна, чем отдельные ценности действий $Q^\pi(s, a)$.  Когда агент находится в безопасном состоянии, выбор действия мало важен; важно само состояние.  Они предложили \textbf{дуэльную архитектуру}\index{dueling network architecture}, которая разделяет финальные слои сети на два потока: один оценивает функцию ценности состояния $V_\theta(s)$, другой --- \textbf{функцию преимущества}\index{advantage function}:
\[
A_\theta(s, a) = Q_\theta(s, a) - V_\theta(s).
\]
Потоки объединяются как:
\begin{equation}
\label{eq:dueling}
Q_\theta(s, a) = V_\theta(s) + \left(A_\theta(s, a) - \frac{1}{|\cA|}\sum_{a'} A_\theta(s, a')\right),
\end{equation}
где вычитание среднего обеспечивает идентифицируемость (поскольку без него $V$ и $A$ не идентифицируемы по отдельности).  Дуэльная архитектура позволяет сети обучаться тому, какие состояния ценны, не обучая воздействию каждого действия, что полезно, когда действия имеют схожий эффект в большинстве состояний.

\subsection{Приоритетное повторное воспроизведение опыта}
\label{subsec:per}

При стандартном повторном воспроизведении переходы извлекаются из $\mathcal{D}$ равномерно.  Шаул и др. (2016) утверждали, что переходы с большими ошибками TD содержат более «удивительную» информацию и должны воспроизводиться чаще.  \textbf{Приоритетное повторное воспроизведение опыта}\index{prioritised experience replay} (PER) назначает каждому переходу $i$ приоритет:
\[
p_i = |\delta_i|^\alpha + \epsilon,
\]
где $\delta_i$ --- наиболее свежая ошибка TD для перехода $i$, $\alpha > 0$ управляет степенью приоритизации, а $\epsilon > 0$ предотвращает нулевой приоритет.  Вероятность выборки: $P(i) = p_i / \sum_j p_j$.  Веса взвешивания по важности $w_i = (N \cdot P(i))^{-\beta}$ (при $\beta \to 1$ по мере продвижения обучения) корректируют неравномерную выборку для сохранения несмещённых оценок градиента.  PER обычно ускоряет обучение и улучшает итоговую производительность, особенно в средах с редкими вознаграждениями, где информативные переходы редки.

% ===========================================================
\section{Анализ сходимости}
\label{sec:convergence-vfa}
% ===========================================================

\subsection{Что можно доказать}

Ландшафт сходимости для аппроксимации функций ценности значительно более неоднозначен, чем для табличного ОП.  В таблице~\ref{tab:convergence-summary} обобщены основные результаты.

\begin{table}[ht]
\centering
\caption{Сводка сходимости для методов аппроксимации функций ценности.  «Сходится» означает сходимость к полезному (возможно, приближённому) решению почти наверное или в среднеквадратичном.  «Может расходиться» означает, что известны контрпримеры сходимости.}
\label{tab:convergence-summary}
\footnotesize
\renewcommand{\arraystretch}{1.2}
\begin{tabular}{@{}llp{2.1cm}p{2.3cm}@{}}
\toprule
\textbf{Метод} & \textbf{Аппрокс.} & \textbf{Выборка} & \textbf{Гарантия} \\
\midrule
МК-градиент & Лин./Нелин. & Любая & Сход. (лок. мин.) \\
Полуград. TD & Линейная & На тек. & Сход. (н.т. TD) \\
Полуград. TD & Линейная & На друг. & Может расход. \\
Полуград. TD & Нелинейн. & На тек. & Может расход. \\
Град. TD (GTD2) & Линейная & На друг. & Сход. (н.т. TD) \\
Полуград. Q-learning & Линейная & На друг. & Может расход. \\
DQN & Нелин. (глуб.) & На друг. & Нет гарантий \\
\bottomrule
\end{tabular}
\end{table}

\subsection{Методы градиентного TD}

Для выхода из смертельной триады при сохранении обучения на другой стратегии несколько авторов разработали \textbf{методы градиентного TD}\index{gradient-TD methods}\,---\,алгоритмы, минимизирующие суррогатную целевую функцию, истинный градиент которой поддаётся вычислению.  Ключевое наблюдение (Саттон и др., 2009) --- определить \emph{проецированную ошибку Беллмана}:
\[
\overline{\mathrm{PBE}}(\mathbf{w}) = \|\Pi_\mu(T^\pi V_\mathbf{w} - V_\mathbf{w})\|_\mu^2,
\]
которая равна нулю в неподвижной точке TD.  В отличие от среднеквадратичной ошибки ценности, градиент $\overline{\mathrm{PBE}}$ можно оценить по одной выборке (с вспомогательным вектором параметров), избегая проблемы двойной выборки.  Полученные алгоритмы (TDC и GTD2) сходятся к $\mathbf{w}_{\mathrm{TD}}$ при выборке на другой стратегии, ломая одну ногу смертельной триады.  Их практическое применение было ограниченным, однако, поскольку они требуют тонкой настройки двух шагов обучения и могут медленно сходиться.

\subsection{Нелинейный случай}

Для глубоких нейронных сетей не известна никакая гарантия сходимости, аналогичная теореме~\ref{thm:linear-td-convergence}, и существуют реальные контрпримеры сходимости с нелинейными аппроксиматорами (Цицикис и Ван Рой, 1997).  Инженерные инновации DQN\,---\,повторное воспроизведение опыта и целевые сети\,---\,являются эвристическими стабилизаторами, хорошо работающими на практике без теоретических гарантий.

Один частичный результат доступен: если потери DQN~\eqref{eq:dqn-loss} рассматривать как задачу обучения с учителем с фиксированными целями (то есть $\theta^-$ постоянно), то стандартное обучение нейронной сети сходится к локальному минимуму $L(\theta)$ для этой фиксированной цели.  Трудность в том, что цель периодически меняется.  Если целевая сеть обновляется бесконечно медленно (формально, $C \to \infty$ при полностью сошедшейся онлайн-сети на каждом шаге), получается итеративная последовательность задач обучения с учителем, неподвижные точки которых связаны с неподвижной точкой Q-learning.  Эта перспектива \emph{подогнанной итерации ценности}\index{fitted value iteration} (Гордон, 1999; Ридмиллер, 2005) обеспечивает некоторое теоретическое обоснование поведения DQN.

% ===========================================================
\section{Приложения и практические соображения}
\label{sec:applications-vfa}
% ===========================================================

\subsection{Инженерия признаков против сквозного обучения}

Центральное проектное решение в любой системе аппроксимации функций ценности --- проектировать ли признаки вручную или обучать их сквозным образом.

\textbf{Признаки, спроектированные вручную} (тайловые кодирования, радиальные базисные функции, базисы Фурье) интерпретируемы, требуют меньше данных и допускают гарантии сходимости при линейном приближении.  Они правильный выбор, когда доменные знания богаты и пространство состояний умеренно сложно.  Например, задача навигации робота на известной карте могла бы использовать в качестве признаков координаты $(x, y)$ робота, его курс и расстояние до ближайшего препятствия\,---\,четырёхмерное представление, сжимающее тысячи данных LIDAR до действительно важных величин.

\textbf{Сквозное глубокое обучение} необходимо, когда сырое наблюдение высокоразмерно (изображения, аудио, текст) и релевантные признаки заранее не известны.  Цена --- значительно большее число параметров, более длительное обучение, сниженная интерпретируемость и более слабые теоретические гарантии.  На практике сквозное глубокое ОП выигрывает от: (i) достаточной ёмкости сети для представления функции ценности; (ii) нормализации признаков для предотвращения взрыва градиентов; (iii) хорошо настроенного расписания $\varepsilon$-жадности для баланса исследования и использования; (iv) достаточно большого буфера повторного воспроизведения для разрыва временных корреляций.

\subsection{Чувствительность к гиперпараметрам}

Алгоритмы глубокого ОП печально известны чувствительностью к гиперпараметрам.  Хендерсон и др. (2018) показали, что заявленная производительность алгоритмов глубокого ОП может варьироваться более чем на порядок в зависимости от случайных начальных состояний, деталей реализации и выбора гиперпараметров.  Следующие гиперпараметры наиболее критически влияют на производительность DQN:
\begin{description}[leftmargin=!, labelwidth=3.5cm]
\item[Размер буфера $N$:] Слишком малый --- буфер заполняется коррелированным недавним опытом; слишком большой --- старые, устаревшие переходы из ранней стратегии преобладают.  Типичные значения: $10^5$--$10^6$.
\item[Частота обновления цели $C$:] Слишком малая --- проблема движущейся цели; слишком большая --- цель сильно отстаёт, замедляя обучение.  Типично: $10^3$--$10^4$ шагов.
\item[Размер пакета $B$:] Большие пакеты дают более стабильные оценки градиента, но увеличивают затраты памяти и вычислений на обновление.  Типично: $32$--$256$.
\item[Расписание $\varepsilon$:] Скорость убывания от $1.0$ до малого финального $\varepsilon$ должна соответствовать задаче; слишком быстрая ведёт к преждевременному использованию; слишком медленная расточительна.
\end{description}

\subsection{Исторические замечания}

Аппроксимация функций в ОП имеет долгую историю.  Сам Беллман признал проклятие размерности ещё в 1950-х годах.  Сходимость линейного TD(0) была установлена Цицикисом и Ван Роем (1997).  Смертельная триада была выявлена и названа Саттоном (1995).  Нейросетевая аппроксимация функций в ОП восходит как минимум к TD-Gammon Тесауро (1992), который научился играть в нарды на уровне, близком к гроссмейстерскому, используя вручную созданную двухслойную сеть, обученную методами TD.  DQN (Мних и др., 2013, в полной версии в 2015) был первой демонстрацией сквозного глубокого ОП, достигшего сверхчеловеческой производительности в разнообразных задачах.  Double DQN (Ван Хасселт и др., 2016), Dueling DQN (Ван и др., 2016) и приоритетное повторное воспроизведение опыта (Шаул и др., 2016) были опубликованы в том же году, совместно образуя «эпоху DQN» в глубоком ОП.  Алгоритм Rainbow (Хессель и др., 2018) объединил все шесть расширений в единого агента, достигнув значительно улучшенной выборочной эффективности.

% ===========================================================
\section*{Упражнения}
\label{sec:exercises-vfa}
\addcontentsline{toc}{section}{Упражнения}
% ===========================================================

\begin{exercise}
\label{ex:curse-dimensionality}
Рассмотрим робота, навигирующего в клеточном мире, описываемом $n$-мерным бинарным вектором состояния (каждое измерение принимает значение $0$ или $1$).
\begin{enumerate}
\item Сколько состояний требует хранения табличная Q-функция с $|\cA| = 4$?
\item Предположим, что вместо этого используется линейная аппроксимация функций с $d = 10n$ признаками, где $d \ll 2^n$. Выразите отношение памяти, требуемой линейным приближением, к табличному представлению как функцию от $n$.
\item При каком значении $n$ линейное приближение впервые требует меньше памяти, чем табличное (при $d = 10n$)?
\item Обсудите одну ошибку приближения, которую вводит линейное представление по сравнению с табличным.
\end{enumerate}
\end{exercise}

\begin{exercise}
\label{ex:projection-compute}
Пусть $|\cS| = 4$, $d = 2$, $\mu = (1/4, 1/4, 1/4, 1/4)$ и
\[
\Phi = \begin{pmatrix} 1 & 0 \\ 1 & 0 \\ 0 & 1 \\ 0 & 1 \end{pmatrix}.
\]
\begin{enumerate}
\item Убедитесь, что $\Phi$ имеет полный ранг по столбцам.
\item Вычислите матрицу проекции $\Pi_\mu = \Phi(\Phi^\top D_\mu \Phi)^{-1} \Phi^\top D_\mu$.
\item Примените проекцию к вектору $v = (3, 1, 2, 4)^\top$.  Каково $\Pi_\mu v$?
\item Проверьте, что $\|v - \Pi_\mu v\|_\mu \leq \|v - f\|_\mu$ для $f = (2, 2, 2, 2)^\top$.
\end{enumerate}
\end{exercise}

\begin{exercise}
\label{ex:semi-gradient-not-gradient}
Рассмотрим среднеквадратичную ошибку ценности $\overline{\mathrm{VE}}(\mathbf{w}) = \sum_s \mu(s)[V^\pi(s) - V_\mathbf{w}(s)]^2$.  Цель TD для перехода $(s, r, s')$: $y = r + \gamma V_\mathbf{w}(s')$.
\begin{enumerate}
\item Запишите $-\frac{1}{2}\nabla_\mathbf{w}[y - V_\mathbf{w}(s)]^2$, рассматривая $y$ как константу (полуградиент).
\item Запишите $-\frac{1}{2}\nabla_\mathbf{w}[y - V_\mathbf{w}(s)]^2$, рассматривая $y$ как функцию от $\mathbf{w}$ (полный градиент).
\item Определите слагаемое, которое полуградиент опускает.  При каком условии это слагаемое равно нулю?
\end{enumerate}
\end{exercise}

\begin{exercise}
\label{ex:td-bound}
Оценка ошибки приближения в неподвижной точке TD (теорема~\ref{thm:td-bound}) гласит: $\|\hat{V} - V^\pi\|_\mu \leq \frac{1}{\sqrt{1-\gamma^2}}\|\hat{V}^* - V^\pi\|_\mu$.
\begin{enumerate}
\item Вычислите коэффициент оценки для $\gamma \in \{0.5, 0.9, 0.99, 0.999\}$.
\item При $\gamma \to 1$ коэффициент расходится.  Дайте интуитивное объяснение, почему бутстрэппинг с дисконтом, близким к единице, усиливает ошибки приближения.
\item Является ли оценка точной в общем случае?  Что должно выполняться в МПР и признаковом наборе для достижения равенства?
\end{enumerate}
\end{exercise}

\begin{exercise}
\label{ex:deadly-triad-identify}
Для каждого из следующих алгоритмов определите, какие элементы смертельной триады присутствуют, и укажите, гарантирована ли сходимость:
\begin{enumerate}
\item Монте-Карло с линейной аппроксимацией функций, на текущей стратегии.
\item Q-learning в табличном МПР с $|\cS| = 10^6$.
\item TD(0) на другой стратегии с нейронной сетью.
\item SARSA с линейной аппроксимацией функций, на текущей стратегии.
\item DQN на Atari.
\end{enumerate}
\end{exercise}
